\section{AI 编译器前端：模型导入、manifest 与 shape verifier}

\subsection{先从一个会悄悄跑错的模型开始}

推理引擎最危险的失败，不是程序立刻崩溃，而是模型能跑、速度也正常、生成结果却已经被污染。假设一个 7B 模型目录里有这些文件：

\begin{lstlisting}[numbers=none,caption={一个本地模型目录的最小文件组}]
config.json
tokenizer.model
tokenizer_config.json
weights-00001.bin
weights-00002.bin
quantization.json
\end{lstlisting}

如果 \code{config.json} 写着 \code{hidden_size = 4096}，但某个 \code{Wq} 权重实际按 \code{[4096, 4096]} 以外的形状保存，后端 linear kernel 仍然可能读到连续字节并给出一个输出。若 tokenizer 的词表大小和 \code{lm_head} 的输出维度差一个特殊 token，采样器也可能只在极少数 prompt 上出错。若某个 int4 权重缺少 group scale，kernel 甚至能按照错误 scale 反量化出一串“看起来像数”的浮点值。

这就是 AI 编译器前端存在的理由。它不只是把文件读进内存，而是要把模型文件转换成一个已检查的 \term{manifest}{清单} 和 typed graph。manifest 不是随手写的说明文档，而是“后端可以相信哪些事实”的合同：有哪些张量、每个张量在哪里、形状是什么、dtype 是什么、量化参数是什么、哪个算子消费它、输出应该是什么。

\begin{keyidea}
模型导入阶段的目标不是“加载成功”，而是建立可验证的语义合同。没有 manifest、shape verifier 和 typed graph，后端 kernel 只能在热路径里猜测输入是否正确。
\end{keyidea}

\subsection{manifest：把文件事实变成编译器事实}

传统编译器的源文件需要词法和语法分析；AI 编译器的模型文件需要 manifest。manifest 至少要收集四类事实：

\begin{itemize}
  \item 模型结构：层数、hidden size、intermediate size、head 数、KV head 数、head dimension、最大上下文长度。
  \item tokenizer 合同：词表大小、特殊 token、BOS/EOS/PAD 规则、是否需要加入前缀空格、byte fallback 或 normalizer。
  \item 权重目录：每个权重张量的名字、文件偏移、字节长度、dtype、shape、逻辑 layout 和 checksum。
  \item 量化元数据：量化格式、group size、scale dtype、zero point 策略、per-tensor/per-channel/per-group 边界。
\end{itemize}

下面是一个压缩后的 manifest 草图。真实工程里可以来自 JSON、二进制索引或模型格式内置 metadata；教材先写成文本，是为了看清不变量。

\begin{lstlisting}[numbers=none,caption={模型 manifest 的核心字段}]
model:
  architecture: decoder_only_transformer
  layers: 32
  hidden_size: 4096
  intermediate_size: 11008
  attention_heads: 32
  kv_heads: 8
  head_dim: 128
  max_context: 4096

tokenizer:
  vocab_size: 32000
  bos_id: 1
  eos_id: 2

tensor:
  name: layers.0.attn.wq.weight
  role: linear_weight
  shape: [4096, 4096]
  dtype: q4_group
  group_size: 64
  scale_dtype: f16
  file: weights-00001.bin
  offset: 1048576
  bytes: 9437184
\end{lstlisting}

manifest 的作用不是增加格式复杂度，而是把“文件里有什么”变成“编译器知道什么”。后续 graph import 不再直接扫文件名；它查询 manifest。后续 shape verifier 不再靠字符串约定；它读 tensor descriptor。后续 backend lowering 不再猜某个权重是不是 int4；它读 dtype 和 quantization descriptor。

\subsection{导入不是字符串匹配}

许多教学代码会用简单规则找权重名，例如看到 \code{wq} 就当作 Q projection，看到 \code{w1} 就当作 MLP gate。真实模型格式会让这种做法很快崩掉：不同模型使用不同命名，某些权重已经被合并为 QKV packed projection，某些模型有 grouped-query attention，某些模型把 \code{lm_head} 和 embedding tied 在一起，某些模型的 RoPE scaling 规则也不同。

前端应该把导入分成三步：

\begin{lstlisting}[numbers=none,caption={模型导入的三阶段}]
parse_files:
  read config, tokenizer config, tensor index, quant metadata
  build raw manifest entries

resolve_schema:
  map model-specific names to canonical roles
  detect tied weights and packed projections
  attach tokenizer and generation contracts

verify_and_import:
  run shape/dtype/quant checks
  produce typed graph nodes and tensor descriptors
\end{lstlisting}

第一步只读事实，不解释太多。第二步把模型方言映射到本引擎的 canonical roles，例如 \code{q_proj.weight}、\code{attention.wq.weight}、\code{layers.0.attention.wq} 都可能被解析为同一个角色。第三步才产生 typed graph。这个顺序很重要：如果一边扫文件一边直接创建 graph node，错误通常会散落在导入代码、运行时代码和后端代码里。

\subsection{schema resolution：把模型方言翻译成统一角色}

不同模型家族最烦人的地方不是数学完全不同，而是同一数学对象有不同名字、不同转置约定和不同融合方式。前端必须把它们翻译成本引擎内部的统一角色。否则每个后端 kernel 都会开始识别模型名字，系统很快失控。

\begin{center}
\begin{tabular}{p{0.28\linewidth}p{0.25\linewidth}p{0.31\linewidth}}
\toprule
外部命名示例 & canonical role & verifier 要确认 \\
\midrule
\code{model.embed_tokens.weight} & token embedding & 行数等于 vocab size，列数等于 hidden \\
\code{self_attn.q_proj.weight} & attention Q & 输出维等于 \code{n_q * head_dim} \\
\code{self_attn.k_proj.weight} & attention K & 输出维等于 \code{n_kv * head_dim} \\
\code{mlp.gate_proj.weight} & MLP gate & 输出维等于 intermediate \\
\code{lm_head.weight} & logits projection & 是否 tied embedding，词表大小一致 \\
\bottomrule
\end{tabular}
\end{center}

有些模型把 Q/K/V 合并成一个大张量，例如 \code{qkv_proj.weight}。这时 schema resolution 不能简单把它当成普通 linear，而要记录切片规则：

\begin{lstlisting}[numbers=none,caption={QKV packed projection 的导入合同}]
PackedQKV:
  source_tensor
  q_range: [0, n_q * head_dim)
  k_range: [q_end, q_end + n_kv * head_dim)
  v_range: [k_end, k_end + n_kv * head_dim)
  storage_layout: row_major | col_major | backend_packed
\end{lstlisting}

如果这个合同不显式保存，后端就可能把 \code{Q/K/V} 的切片顺序写死在 kernel 里。等另一个模型把顺序改成 \code{Q/V/K}，程序仍然能跑，但 attention 会稳定地错。

\subsection{一层 decoder 的 shape 推导例子}

把 shape verifier 写成代码之前，要能在纸上推导一层。设配置为：

\begin{lstlisting}[numbers=none,caption={一组 7B 级 decoder 配置}]
hidden_size = 4096
intermediate_size = 11008
attention_heads = 32
kv_heads = 8
head_dim = 128
\end{lstlisting}

则 verifier 应该推出：

\begin{center}
\begin{tabular}{p{0.22\linewidth}p{0.28\linewidth}p{0.34\linewidth}}
\toprule
张量角色 & 期望形状 & 推导来源 \\
\midrule
Q weight & \code{[4096,4096]} & \code{n_q * d = 32 * 128} \\
K weight & \code{[1024,4096]} & \code{n_kv * d = 8 * 128} \\
V weight & \code{[1024,4096]} & 同 K \\
O weight & \code{[4096,4096]} & query heads 拼回 hidden \\
MLP gate/up & \code{[11008,4096]} & intermediate by hidden \\
MLP down & \code{[4096,11008]} & hidden by intermediate \\
RMSNorm gamma & \code{[4096]} & 每个 hidden 维度一个权重 \\
\bottomrule
\end{tabular}
\end{center}

这个例子很重要，因为它能立即抓住 GQA 配置错误。若 K/V 权重被导入成 \code{[4096,4096]}，但配置写着 \code{kv_heads=8}，前端必须拒绝或明确标记这是非 GQA 模型；不能让 attention kernel 到运行时才发现 \code{kv_head} 索引和权重输出维度对不上。

\subsection{shape verifier：让每条边都有可证明的形状}

shape verifier 是 AI 编译器前端的语义分析器。它的输入是 manifest 和 graph schema，输出是带 shape 的 typed graph，或者一组明确诊断。以一层 decoder block 为例，若 hidden size 记为 \code{H}，head 数为 \code{A}，KV head 数为 \code{KVA}，head dimension 为 \code{D}，那么至少有：

\begin{lstlisting}[numbers=none,caption={decoder 层 shape 约束}]
H == A * D
K_shape == [KVA * D, H]
V_shape == [KVA * D, H]
Q_shape == [A * D, H]
Wo_shape == [H, A * D]
Wgate_shape == [I, H]
Wup_shape == [I, H]
Wdown_shape == [H, I]
RMSNorm_gamma_shape == [H]
KV_cache_shape == [layers, batch, KVA, max_context, D]
\end{lstlisting}

这些约束决定的不只是正确性，还决定性能路径。若 \code{KVA < A}，这是 grouped-query attention，decode attention kernel 要知道多个 query head 共享一组 K/V head；若 \code{D} 不是 SIMD-friendly 的倍数，CPU 后端要准备尾部处理；若 \code{I} 和 tile size 不对齐，MLP kernel 需要 remainder path。shape verifier 把这些事实提前暴露出来，而不是让后端在热循环里临时发现。

\begin{systemdiagram}{manifest 到 typed graph 的前端流水线}
\begin{pdfdiagram}[x=1cm,y=1cm]
  \node[book accent node,text width=2.35cm] (files) at (0,1.45) {模型文件\\config/weights};
  \node[book teal node,text width=2.35cm] (manifest) at (3.15,1.45) {manifest\\张量清单};
  \node[book purple node,text width=2.35cm] (schema) at (6.3,2.35) {schema\\角色映射};
  \node[book amber node,text width=2.35cm] (verify) at (6.3,0.55) {verifier\\shape\\dtype/quant};
  \node[book navy node,text width=2.35cm] (graph) at (9.45,1.45) {typed graph\\后端输入};

  \draw[book strong arrow] (files) -- (manifest);
  \draw[book strong arrow] (manifest) -- (schema);
  \draw[book strong arrow] (manifest) -- (verify);
  \draw[book weak arrow] (schema) -- (verify);
  \draw[book strong arrow] (verify) -- (graph);

  \node[book label,text width=9.8cm] at (4.75,-0.95) {前端先把文件事实稳定成 manifest，再把模型方言解析成角色，最后由 verifier 产出可被后端相信的 typed graph。};
\end{pdfdiagram}
\diagramnote{本图要证明的命题是：模型导入不是直接读权重跑 kernel，而是 manifest、schema 和 verifier 共同建立编译器事实。}
\end{systemdiagram}

\subsection{dtype verifier：不是每个字节数组都能进 kernel}

shape 对了，dtype 仍然可能错。一个 linear 权重可能是 fp16、bf16、int8、int4 group-wise，也可能是某个后端专用 packed format。dtype verifier 要检查三层事实：

\begin{itemize}
  \item 逻辑 dtype：算子语义中权重代表什么数值类型，例如 f16 weight 或 q4 group weight。
  \item 存储 dtype：文件中每个元素或压缩块怎样编码，例如 4-bit pair、int8、f16 scale。
  \item 计算 dtype：kernel 中使用什么累加器和输出类型，例如 fp32 accumulator、int32 dot accumulator、f16 output。
\end{itemize}

这三层不能混为一谈。int4 权重的逻辑值是近似浮点权重，存储上可能两个 4-bit value 打包在一个字节，计算时可能先反量化到寄存器再做 FMA，也可能走整数 dot 再 requantize。若 manifest 只写 \code{dtype=int4}，后端还不知道 group size、scale layout、zero point、padding 和尾部怎么处理。

\begin{lstlisting}[numbers=none,caption={量化权重 verifier 要检查的边界}]
for each quantized weight tensor:
  require group_size > 0
  require input_channels % group_size == 0 or tail_policy exists
  require scale tensor exists
  require scale length matches groups
  require packed byte length matches shape and bit width
  require backend supports this quant format
\end{lstlisting}

这些检查看似繁琐，但比在 kernel 中读越界 scale 或错误解码权重便宜得多。前端拒绝模型是清晰错误；后端悄悄算错 logits 是高成本错误。

\subsection{verifier 负例矩阵}

前端测试不能只加载一个正确模型。必须构造坏 manifest，证明 verifier 会在正确层次拒绝它。最小负例矩阵如下：

\begin{center}
\begin{tabular}{p{0.24\linewidth}p{0.36\linewidth}p{0.24\linewidth}}
\toprule
负例 & 应报位置 & 典型风险 \\
\midrule
embedding 行数少于 vocab & tokenizer/embedding verifier & token id 越界 \\
K/V head 数与权重维度不一致 & shape verifier & GQA attention 读错 head \\
QKV packed 切片顺序缺失 & schema verifier & Q/K/V 语义交换 \\
int4 scale 数量不足 & quant verifier & kernel 越界或误反量化 \\
group size 为 0 或不支持 & quant verifier & 除零或 unsupported plan \\
RoPE head dim 为奇数 & position verifier & 二维旋转对不完整 \\
max context 小于 prompt & runtime admission & KV cache 越界 \\
GPU 不支持目标 dtype & backend capability & 静默 fallback 或失败 \\
\bottomrule
\end{tabular}
\end{center}

每个负例都应该检查诊断内容，而不只是“失败了”。比如 K/V head 维度错误的诊断要写出 \code{kv_heads * head_dim} 的推导；int4 scale 数量错误要写出 expected scale count；GPU dtype 不支持要写出 requested profile 和 backend capability。

\subsection{RoPE 和位置参数也要验证}

RoPE 参数如果错，shape 仍然可能全部正确。前端至少要检查：

\begin{itemize}
  \item \code{head_dim} 是否能拆成二维旋转对，通常要求偶数。
  \item RoPE base、scaling 类型、原始最大位置和扩展规则是否存在。
  \item prefill 与 decode 使用同一 position 语义。
  \item sliding window 或 prefix cache 是否记录 base position。
\end{itemize}

\begin{lstlisting}[numbers=none,caption={RoPE verifier 字段}]
RopeDesc:
  head_dim
  base_theta
  scaling_policy
  original_max_position
  effective_max_position
  position_base_policy
\end{lstlisting}

如果模型使用扩展上下文策略，而 manifest 没记录 scaling policy，引擎可以选择拒绝，而不是默认用基础 RoPE。默认继续跑会得到 shape 正确但长上下文错误的结果。

\subsection{tokenizer 也是编译合同的一部分}

很多人把 tokenizer 当作预处理，但本地推理引擎不能这样切开。tokenizer 决定 token id 序列，token id 决定 embedding 访问和位置编码，最终影响每一步 logits。若 tokenizer 的 BOS/EOS 规则、normalizer 或特殊 token 与模型训练时不一致，后端再快也只是高效地产生错误输入。

前端至少要验证：

\begin{itemize}
  \item tokenizer 词表大小等于 embedding 行数，或有明确扩展规则。
  \item \code{bos_id}、\code{eos_id}、\code{unk_id} 等特殊 token 在词表范围内。
  \item generation config 中的 stop token 与 tokenizer 合同一致。
  \item prompt encode 后的 token 数不会超过当前 session 的 max context。
  \item tokenizer 输出 dtype 能被 embedding kernel 接受，通常是 32 位整数索引。
\end{itemize}

这部分和编译器的符号解析有相似性。源代码里的名字必须解析到某个声明；prompt 里的文本必须解析到模型词表中的 token id。若解析规则错了，后面所有 IR 和 kernel 都建立在错误输入上。

\subsection{诊断：错误要停在前端，而不是后端崩溃}

shape verifier 的诊断要像编译器错误，而不是像运行时崩溃。一个有用诊断应该包含：哪个 tensor、期望形状、实际形状、这个期望来自哪个 config 字段或上游算子，以及可能的修复方向。

\begin{lstlisting}[numbers=none,caption={好的模型导入诊断示例}]
error: tensor layers.12.attn.wk.weight has incompatible shape
  expected: [1024, 4096]
    reason: kv_heads(8) * head_dim(128), hidden_size(4096)
  actual:   [4096, 4096]
  note: this looks like full multi-head K projection,
        but config declares grouped-query attention
\end{lstlisting}

这类诊断让使用者知道是模型文件、config、转换脚本还是引擎 schema 出了问题。反过来，如果只在后端报 \code{invalid argument} 或 \code{segmentation fault}，调试会被迫从机器码和二进制文件倒推语义，成本完全不对等。

\subsection{前端产物：typed graph 必须足够给后端使用}

前端通过后，输出不应该只是“一个 graph 指针”。它至少要包含：

\begin{itemize}
  \item Graph nodes：算子类型、输入输出边、状态边和执行阶段标签。
  \item Tensor descriptors：shape、stride、dtype、layout、quantization、file location、ownership。
  \item Symbol table：canonical tensor name 到 descriptor 的映射。
  \item State descriptors：KV cache、position、session buffer 和 sampler state。
  \item Verification report：通过了哪些约束，哪些路径需要 fallback 或 tail handling。
\end{itemize}

这些产物会直接进入第 6 章讲过的后续流水线：Graph IR rewrite、Tensor IR layout 选择、memory planning、CPU/GPU lowering 和 runtime dispatch。前端如果偷懒，后端就会被迫补课；而后端补课通常发生在最难调试、最要求性能稳定的地方。

\subsection{真实模型资产闭环：从目录到第一条 token trace}

一个本地推理引擎不能只验证 toy fixture。真正有求职说服力的前端报告，应能把一个真实模型目录推进到第一条可解释 token trace。这里的“真实”不是必须一开始支持所有格式，而是至少支持一种明确资产形态，并把文件事实、tokenizer、权重索引、typed graph 和 reference decoder 连接起来。

\begin{lstlisting}[numbers=none,caption={真实模型导入的最小验收路径}]
model_asset
  -> parse metadata and tokenizer contract
  -> build tensor index with offsets, byte lengths, dtype
  -> verify shapes, quant descriptors, checksums
  -> create typed graph and state descriptors
  -> mmap or load weights under ownership rules
  -> run tokenizer on fixed prompt
  -> run one reference decode step
  -> dump checkpoint trace and logits checksum
\end{lstlisting}

这条路径比“模型能加载”更严格。它要求每个阶段都有可保存的产物：\filepath{manifest.json}、\filepath{tensor-index.csv}、\filepath{verify-report.txt}、\filepath{graph-ir.txt}、\filepath{trace.csv} 和 \filepath{logits-golden.txt}。后端优化可以晚一点做，但前端资产闭环必须先稳住，否则后面的 kernel benchmark 很可能在错误输入上跑得很快。

\subsection{GGUF 和 safetensors：两类不同边界}

真实模型常见资产形态至少有两类。GGUF 更像一个把 metadata、tensor directory 和权重字节放在同一容器里的本地推理资产；safetensors 更强调安全的张量存储和 metadata，tokenizer/config 通常以旁路 JSON 文件存在。教材不要求一开始同时实现两者，但要把它们的工程边界讲清楚。

\begin{center}
\small
\begin{tabular}{p{0.20\linewidth}p{0.34\linewidth}p{0.32\linewidth}}
\toprule
资产形态 & 前端要提取 & 常见错误边界 \\
\midrule
GGUF & metadata、tensor name、offset、dtype、shape、quant type & metadata 方言、量化块解释、tokenizer 字段不完整 \\
safetensors & JSON header、tensor offsets、dtype、shape、metadata & config/tokenizer 分离，权重名方言需要 schema resolution \\
tokenizer JSON/model & vocab、normalizer、pre-tokenizer、special tokens & BOS/EOS 规则、byte fallback、词表和 embedding 不一致 \\
generation config & stop token、max length、sampling defaults & sampler 和 tokenizer 合同不一致 \\
\bottomrule
\end{tabular}
\end{center}

这张表的结论是：模型格式不是普通文件读取问题。GGUF 的 tensor directory 能帮你快速建立权重索引，但你仍要解释量化块；safetensors 的张量边界清晰，但模型结构和 tokenizer 合同往往要从其他文件合并。无论哪种格式，前端最终都必须产出同一种内部 manifest，否则后端会被格式细节污染。

\subsection{tensor index：mmap 之前先证明字节范围}

权重文件最好不要一开始就读成一堆 \code{std::vector<float>}。真实模型可能十几 GB，前端需要先建立 tensor index：每个张量的文件、偏移、字节长度、对齐、dtype、shape、checksum 和逻辑角色。只有 index 通过验证后，runtime 才能决定 mmap、按需加载、预取或复制到 backend packed cache。

\begin{lstlisting}[numbers=none,caption={tensor index 的最低字段}]
TensorIndexEntry:
  canonical_name
  source_name
  file_path
  offset_bytes
  length_bytes
  alignment
  shape
  storage_dtype
  logical_dtype
  quant_descriptor_id
  checksum
  mmap_allowed
\end{lstlisting}

验证规则要非常机械。张量字节范围不能越过文件大小；两个张量范围不能意外重叠，除非格式明确允许 alias 或 tied weight；offset 要满足格式和后端对齐要求；packed quant tensor 的字节长度要能由 shape、bit width、block size 和 scale metadata 推导出来。若这些规则不通过，错误必须停在导入阶段。

\begin{lstlisting}[numbers=none,caption={字节范围 verifier 的拒绝条件}]
reject if offset + length > file_size
reject if range overlaps another tensor without alias record
reject if length != expected_storage_bytes(shape, dtype, quant)
reject if checksum exists and does not match
reject if backend requires alignment and offset violates it
\end{lstlisting}

这一步能阻止很多“后端崩溃”。例如 int4 权重少了最后一个 scale block，kernel 可能只在某个 output tile 读到越界；如果 tensor index 阶段已经计算 expected bytes，错误会变成一条清楚诊断，而不是随机 logits 漂移。

\subsection{tokenizer golden：真实 prompt 的第一道 oracle}

真实模型链路的第一个 golden output 不应该是最终生成文本，而应该是 tokenizer 输出。固定一个短 prompt，例如 \code{"Hello"} 或中文短句，保存 token id 序列、BOS/EOS 插入规则和特殊 token 处理。若 tokenizer 都不一致，后续 embedding、RoPE、KV cache 和 logits 全部没有比较意义。

\begin{lstlisting}[numbers=none,caption={tokenizer golden 报告}]
tokenizer_golden:
  model_id
  tokenizer_hash
  prompt_utf8_hex
  add_bos: true
  add_eos: false
  tokens: [1, 15043]
  token_count: 2
  vocab_size: 32000
  embedding_rows: 32000
\end{lstlisting}

报告里保存 \code{prompt\_utf8\_hex} 是为了避免不可见空格、换行、全角字符、Unicode normalization 和编辑器编码把实验搞乱。中文 prompt 尤其要保存原始 UTF-8 字节，因为不同 tokenizer 对空格、字节 fallback 和 normalization 的处理差异会直接改变 token id。

\subsection{第一条真实 trace：不要直接对最终文本}

前端闭环的第一条 trace 可以只跑一层或一个 decode step，但必须包含中间 checkpoint。建议固定字段如下：

\begin{lstlisting}[numbers=none,caption={真实模型 first-token trace schema}]
trace_row:
  stage
  layer
  position
  tensor_name
  shape
  stride
  dtype
  layout
  checksum
  max_abs
  max_rel_error_vs_reference
\end{lstlisting}

最小 checkpoint 顺序为：

\begin{lstlisting}[numbers=none,caption={first-token trace checkpoint}]
tokenizer
embedding
layer0.input_rmsnorm
layer0.qkv_linear
layer0.rope_q
layer0.rope_k
layer0.kv_append
layer0.attention_scores
layer0.attention_out
layer0.mlp_gate_up
layer0.block_out
final_norm
lm_head_logits
sampler_selected_token
\end{lstlisting}

如果项目暂时没有完整真实模型后端，也可以先只执行 reference path 和 shape trace。关键是产物要固定下来：同一模型资产、同一 prompt、同一 commit 应输出同一批 checkpoint checksum。后续 SIMD kernel、paged KV、int4 weight 或 GPU backend 只要在某一步开始偏离，就能定位到具体阶段。

\subsection{mmap 权重加载的所有权合同}

真实权重导入通常会使用 mmap 或分块读取。mmap 的好处是避免一次性复制整个文件，并允许按需触页；代价是生命周期、page fault、文件替换和错误边界更复杂。前端 manifest 必须把“张量视图”和“文件映射”分开建模。

\begin{lstlisting}[numbers=none,caption={mmap 权重所有权状态机}]
Mapped:
  file descriptor validated
  mapping created
  tensor index ranges point into mapping

PinnedForPlan:
  executable plan holds references to tensor views
  file identity and checksum recorded

Packed:
  backend-specific packed weights created
  packed cache has its own ownership

Released:
  no live tensor view may point into unmapped bytes
\end{lstlisting}

错误案例很具体：如果 runtime 在 packing 后释放 mmap，但某个 fallback kernel 仍指向原始 weight view，短测试可能不触发，长请求或错误路径才崩溃。解决办法不是“别释放”，而是让 executable plan 明确记录每个 segment 使用原始 view、packed cache 还是 device buffer。所有权状态必须能被 verifier 检查。

\subsection{真实模型导入的负例包}

toy 模型负例不够。真实资产导入要构造能模拟转换器错误的负例包：

\begin{center}
\small
\begin{tabular}{p{0.26\linewidth}p{0.35\linewidth}p{0.25\linewidth}}
\toprule
负例 & 应该被谁拒绝 & 若漏掉会怎样 \\
\midrule
tokenizer vocab 比 embedding 多 1 & tokenizer/embedding verifier & 特殊 token 偶发越界 \\
QKV packed 顺序写错 & schema resolver & attention 语义交换但形状正确 \\
tensor offset 少 16 字节对齐 & tensor index verifier & SIMD load 或 mmap view 边界异常 \\
q4 scale block 少一个 & quant verifier & 某个 tile 反量化读错 \\
RoPE scaling metadata 缺失 & position verifier & 短上下文通过，长上下文漂移 \\
lm head tied weight 未声明 alias & alias verifier & memory planner 误释放或重复加载 \\
config max context 大于 KV plan 容量 & admission/runtime verifier & decode 中途越界或拒绝太晚 \\
\bottomrule
\end{tabular}
\end{center}

这些负例训练的是“错误停在哪一层”。如果 QKV 顺序错被 attention kernel 才发现，说明 schema 解析太弱；如果 vocab mismatch 到 sampler 才发现，说明 tokenizer 合同没有前置；如果 q4 scale 缺失到 AVX2 kernel 才崩，说明 quant verifier 没有完成字节账本。

\subsection{导入报告和求职证据}

第三册项目的导入报告应当可以直接放进作品集。它不需要夸大支持所有模型，但必须诚实且可复现：

\begin{lstlisting}[numbers=none,caption={真实模型导入报告目录}]
model-import-report/
  asset-summary.txt
  commands.txt
  environment.txt
  manifest.json
  tensor-index.csv
  verifier-report.txt
  tokenizer-golden.txt
  graph-ir.txt
  first-token-trace.csv
  logits-checksum.txt
  unsupported-features.txt
\end{lstlisting}

\filepath{unsupported-features.txt} 很重要。若当前只支持 fp16 safetensors，不支持 GGUF q4；只支持独立 Q/K/V，不支持 packed QKV；只支持基础 RoPE，不支持某种 scaling；都应该明确写出来。生产级工程不是永远支持一切，而是知道自己支持什么、拒绝什么、如何诊断。

\subsection{真实 LLM 工业链条：从资产到第一条 token trace}

第三册不能只停在“manifest 应该有什么”。真实 LLM 工业链条至少跨过 tokenizer、权重格式、metadata、mmap、shape verifier、graph import、memory plan、backend plan、KV cache 和 sampler。任何一环靠猜，后面的 kernel 再快也只是把错误更快地扩散。

\begin{center}
\small
\begin{tabular}{p{0.21\linewidth}p{0.33\linewidth}p{0.30\linewidth}}
\toprule
链条阶段 & 必须建立的事实 & 典型工业格式或系统 \\
\midrule
asset discovery & 文件集合、分片、checksum、版本 & safetensors index、GGUF metadata \\
tokenizer import & vocab、normalizer、special token、byte fallback & tokenizer.json、SentencePiece、GGUF tokenizer \\
tensor index & name、offset、bytes、dtype、shape、checksum & safetensors header、GGUF tensor table \\
schema resolution & 外部名字映射到 canonical role & Llama/Mistral/Qwen 等方言 \\
quant descriptor & group size、scale、zero point、packing order & GGUF q4/q5/q8、GPTQ/AWQ 产物 \\
typed graph & 每个 op 的 shape/dtype/layout/state edge & 本书 Graph IR、ONNX/MLIR 类系统 \\
memory plan & activation、workspace、KV、packed cache & arena、mmap、device buffer \\
backend plan & CPU/GPU kernel、ISA、fallback、packing & ggml/llama.cpp、oneDNN、TensorRT-LLM \\
first token trace & 每层 checkpoint、logits、top-k、sampler & oracle CSV、profiler event \\
\bottomrule
\end{tabular}
\end{center}

这张表给了工程边界。safetensors 的优势是头部记录 tensor 名、dtype、shape 和 data offsets，适合安全地 mmap 和校验范围；GGUF 的优势是把模型 metadata、tokenizer、量化格式和 tensor table 放在一个文件里，适合本地推理分发。无论选择哪一种，本引擎内部都不应该让后端直接依赖外部文件格式。外部格式先转换成 manifest 和 tensor descriptor，后端只消费已验证合同。

\subsection{工业导入的负载顺序}

真实导入流程应按“先轻后重”组织，避免读了几十 GB 权重后才发现 tokenizer 或 shape 不匹配。

\begin{lstlisting}[numbers=none,caption={真实模型导入的推荐顺序}]
1. read small metadata:
   config, tokenizer metadata, tensor index, quant metadata

2. verify global contracts:
   vocab_size, hidden_size, layers, heads, kv_heads,
   head_dim, rope config, max_context, tied weights

3. verify tensor table:
   names, canonical roles, offsets, byte ranges,
   dtype, shape, quant scale ranges, checksums

4. build typed graph:
   op nodes, tensor descriptors, state edges,
   tokenizer contract, generation contract

5. build memory and backend plan:
   mmap views, packed weight cache, KV cache budget,
   CPU/GPU kernel selection, fallback rules

6. run first-token trace:
   tokenizer ids, embedding, each decoder checkpoint,
   logits checksum, top-k, sampled token
\end{lstlisting}

这个顺序的关键是：metadata 和 tensor index 先闭合，权重字节后进入热路径。若 \code{lm\_head} 和 tokenizer vocab 不一致，应在读取 tensor table 时拒绝；若 RoPE scaling metadata 缺失，应在 graph import 时拒绝；若 q4 scale 数量不对，应在 quant descriptor 验证时拒绝；不能等到 AVX2 kernel 读错 scale 后才崩。

\subsection{GGUF 和 safetensors 的导入差异}

两类格式的工程取舍不同。safetensors 更像“安全 tensor 容器”：它强调 header、shape、dtype 和 data offsets，不执行任意代码，适合 HuggingFace 生态权重分片。GGUF 更像“本地推理包”：它把模型 metadata、tokenizer、量化类型和 tensor table 一起打包，适合 llama.cpp/ggml 风格的本地加载。

\begin{center}
\small
\begin{tabular}{p{0.18\linewidth}p{0.33\linewidth}p{0.33\linewidth}}
\toprule
维度 & safetensors & GGUF \\
\midrule
metadata & 通常还要结合 config/tokenizer 文件 & 文件内包含大量模型和 tokenizer metadata \\
tensor table & header 中有 name/shape/dtype/offset & header 和 tensor info 记录 name/type/offset \\
量化 & 常见 fp16/bf16 或外部量化产物 & 内建多种 ggml quant type \\
tokenizer & 通常来自 tokenizer.json 或 tokenizer.model & 可随文件存储 tokenizer 字段 \\
mmap & 可按 data offsets 建 view & 可按 tensor offset 建 view，但需处理 alignment \\
主要风险 & 多文件版本不一致、tokenizer/config 漂移 & 量化类型解释、metadata 兼容、版本差异 \\
\bottomrule
\end{tabular}
\end{center}

本书项目可以先支持一种格式，但报告必须写明支持边界。只支持 safetensors fp16，并不丢人；声称支持 GGUF 却不解释 q4 block layout、scale 元数据和 tokenizer 字段，才是危险。工业链条的第一原则是可拒绝、可诊断、可复现。

\subsection{第一条 token trace 是导入闭环的验收}

导入成功不能只看“模型文件打开了”。最小验收是一条固定 prompt 的 first-token trace。它应该记录从 prompt 到 sampled token 的每个关键边界：

\begin{lstlisting}[numbers=none,caption={first-token trace 的核心字段}]
FirstTokenTrace:
  model_asset_id
  tokenizer_id
  prompt_text
  token_ids
  graph_digest
  memory_plan_digest
  backend_plan_id
  checkpoints:
    embedding_checksum
    layer0_rmsnorm_checksum
    layer0_qkv_checksum
    layer0_rope_checksum
    layer0_attention_checksum
    layer0_mlp_checksum
    final_logits_checksum
  topk:
    ids
    values
  sampled_token
  tolerance_profile
\end{lstlisting}

这条 trace 同时服务正确性和求职作品集。它能证明 tokenizer、模型 metadata、张量 shape、layout、KV state、后端计划和 sampler 确实连上了。若以后换量化格式、换 packed layout 或换 GPU backend，trace 能指出漂移从哪一层开始。

\subsection{本节验收线}

读完本节，应该能说明 AI 编译器前端在本地推理引擎中的职责：

\begin{itemize}
  \item manifest 把模型目录、权重文件、tokenizer 和量化 metadata 转换成编译器事实。
  \item schema resolution 把不同模型的命名方言映射到统一算子角色。
  \item shape verifier 用 config 和算子合同推导每条边的形状，不让后端猜。
  \item dtype 和 quant verifier 区分逻辑 dtype、存储 dtype 和计算 dtype。
  \item tokenizer 合同必须和 embedding、generation config、max context 一起检查。
  \item 好诊断应该指出 tensor、期望、实际、推导来源和可能错因。
  \item 真实模型导入要产出 manifest、tensor index、tokenizer golden、typed graph 和 first-token trace。
  \item mmap 权重视图、packed weight cache 和 device buffer 要有明确所有权状态。
  \item 负例包要证明错误停在前端，而不是落到后端 kernel 或 sampler 才暴露。
  \item 真实 LLM 工业链条要从外部资产导入到 first-token trace，而不是只加载一个二进制权重数组。
\end{itemize}

下一步进入 runtime 和 memory planner 时，这些前端产物会成为所有优化的前提：只有 typed graph 的边界稳定，内存生命周期、kernel selection 和 CPU/GPU 后端计划才有可信输入。
